\newpage\EMPHASIZE{32. Enumerations}

\textsc{Objectives}
\begin{itemize}
\item Create enums
\item Use enums
\item Understand the relationship between enums and constants
\end{itemize}

Up to this point we have been using types provided by C++. Using these
types we have also been using arrays and pointer types. Now we will
create our own types using \textit{enum}. The \textit{enum} basically
creates a type made up of integer constants.


\newpage\EMPHASIZE{Enum}

There's another way to rewrite a previous example:

\begin{consolethree}[escapeinside=||]
typedef int EmployeeCode;

const EmployeeCode CEO = 0;
const EmployeeCode MANAGER = 1;
const EmployeeCode FULLTIME = 2;
const EmployeeCode PARTTIME = 3;

EmployeeCode code = CEO;
\end{consolethree}

Here's how you do it:

\begin{consolethree}[escapeinside=||]
#include <iostream>

int main()
{
    enum EmployeeCode
    {
        CEO, MANAGER, FULLTIME, PARTTIME
    };
    EmployeeCode code = CEO;
    std::cout << CEO << std::endl;
    std::cout << code << std::endl;
    return 0;
}
\end{consolethree}

Here comes some \textit{enum} jargon ...

\begin{consolethree}[escapeinside=||]
enum EmployeeCode
{
    CEO, MANAGER, FULLTIME, PARTTIME
};
\end{consolethree}

The format of creating an enumeration is

\begin{consolethree}[escapeinside=||]
enum [enumeration name]
{
    [enumerator-1], ...,
    [enumerator-n]
};
\end{consolethree}

\begin{ex}
Declare a \textit{Grade} enumeration containing
enumerators \textit{A}, \textit{B}, \textit{C}, \textit{D}, and \textit{F}.
Declare variable \textit{johnDoeGrade} of type \textit{Grade} and initialize
it with \textit{A}. Print \textit{johnDoeGrade} and then print \textit{A}.
\end{ex}


\newpage\EMPHASIZE{Enum as a type}

So what's the difference between

\begin{consolethree}[escapeinside=||]
typedef int EmployeeCode;

const EmployeeCode CEO = 0;
const EmployeeCode MANAGER = 1;
const EmployeeCode FULLTIME = 2;
const EmployeeCode PARTTIME = 3;
\end{consolethree}

and

\begin{consolethree}[escapeinside=||]
enum EmployeeCode
{
    CEO, MANAGER, FULLTIME, PARTTIME
};
\end{consolethree}

(besides the fact that the second produces a shorter code)???

The enumeration \textit{EmployeeCode} is actually \textbf{a new
type}. This means that you can actually have this:

\begin{ex}
Try compiling this program:

\begin{consolethree}[escapeinside=||]
enum EmployeeCode
{
    CEO, MANAGER, FULLTIME, PARTTIME
};

void f(EmployeeCode c)
{}

void f(int i)
{}

int main()
{
    return 0;
}
\end{consolethree}
\end{ex}


\newpage\EMPHASIZE{Enumerators and Integers}

You might say ... ``Hang on there! What do you mean that the
\textit{EmployeeCode} is a new type??? When I print an \textit{EmployeeCode}
variable like this ...''

\begin{consolethree}[escapeinside=||]
enum EmployeeCode
{
    CEO, MANAGER, FULLTIME, PARTTIME
};

std::cout << CEO << std::endl;
\end{consolethree}

``I see that \textit{CEO} \textit{is} 0!!! So it \textit{is} an integer!!!
Isn't it???''

Yes and no.

There's actually an automatic \textit{typecasting} between
enumerators and integers. The statements:

\begin{consolethree}[escapeinside=||]
std::cout << CEO << std::endl;
std::cout << MANAGER << std::endl;
\end{consolethree}

are really the same as

\begin{consolethree}[escapeinside=||]
std::cout << int(CEO) << std::endl;
std::cout << int(MANAGER) << std::endl;
\end{consolethree}

This means that \textit{CEO}, because it's the
\textbf{first} enumerator of an enumeration, is automatically
typecast to the integer \textbf{0} if necessary, \textit{MANAGER}
is automatically typecast to the integer 1 if necessary, etc.

That example shows you that \textit{EmployeeCode} enumerator is
automatically typecast to an \textit{int} when necessary. What about from
an \textit{int} to \textit{EmployeeCode} enumerator? Try this:

\begin{consolethree}[escapeinside=||]
enum EmployeeCode
{
    CEO, MANAGER, FULLTIME, PARTTIME
};

EmployeeCode code = EmployeeCode(1);
\end{consolethree}

The last statement is the same as:

\begin{consolethree}[escapeinside=||]
EmployeeCode code = MANAGER;
\end{consolethree}

Now let's see if the typecast is also automatic. Try this:

\begin{consolethree}[escapeinside=||]
enum EmployeeCode
{
    CEO, MANAGER, FULLTIME, PARTTIME
};

EmployeeCode code = 1; // 1 becomes EmployeeCode(1)??
\end{consolethree}

The point: An integer is \textbf{not} automatically typecast to
an enumerator.

Since C/C++ will automatically typecast enumerators to integers if
necessary, we can use it as though it's an integer. Here
are some examples:

Run this:

\begin{consolethree}[escapeinside=||]
enum operation {ADD, SUB, MUL, DIV};

operation op = ADD;
int x = 10, y = 20;

switch (op)
{
    case ADD: std::cout << x + y; break;
    case SUB: std::cout << x - y; break;
    case MUL: std::cout << x * y; break;
    case DIV: std::cout << x / y; break;
}
\end{consolethree}

Try changing the \textit{op} to different values.

\begin{ex}
Is this a valid code segment? If so, what is the
output? If it isn't explain why.

\begin{consolethree}[escapeinside=||]
enum operation {ADD, SUB, MUL, DIV};

std::cout << ADD + SUB + MUL + DIV << std::endl;
\end{consolethree}
\end{ex}


\newpage\EMPHASIZE{The integer value of an enumerator}

You can actually change the integer value associated with an enumerator.
Try this:

\begin{consolethree}[escapeinside=||]
enum CutOff
{
    A = 90, B = 80, C = 70, D = 60, F = 0
};

std::cout << B << std::endl;
CutOff x = B;
std::cout << x << std::endl;
\end{consolethree}

Now try this:

\begin{consolethree}[escapeinside=||]
enum X
{
    a = 5, b, c, d = 42, e, f
};

std::cout << a << ' ' << b << ' ' << c << ' '
          << d << ' ' << e << ' ' << f << std::endl;
\end{consolethree}

And ... you \textbf{have} to try this because
it's a common gotcha:

\begin{consolethree}[escapeinside=||]
enum X {a = 1.1};
\end{consolethree}

The point is that enumerators can only be initialized with integer
values.

\begin{ex}
Either write down the output of the following code segment or
find the error:

\begin{consolethree}[escapeinside=||]
enum X
{
    a, b, c, d = 42, e, f
};

X foo = b;
X bar = f;
std::cout << b << f << std::endl;
\end{consolethree}
\end{ex}

\begin{ex}
Either write down the output of the following code segment or
find the error:

\begin{consolethree}[escapeinside=||]
enum Y
{
    a = 1, b, c, d = 4.5, e = 5.6, f = 6.7
};

X foo = b;
std::cout << b << std::endl;
\end{consolethree}
\end{ex}


\newpage\EMPHASIZE{Enum Without Enumeration Name}

Notice that

\begin{consolethree}[escapeinside=||]
enum EmployeeCode
{
    CEO, MANAGER, FULLTIME, PARTTIME
};
\end{consolethree}

actually does \textbf{two} things:

\begin{itemize}
\item It creates four constants (enumerators)
\item It gives a type (enumeration) for the four constants.
\end{itemize}

You can actually use the \textit{enum} to create constants without
creating an enumeration. Try this:

\begin{consolethree}[escapeinside=||]
#include <iostream>

int main()
{
    enum {A, B, C, D, E};
    std::cout << A << std::endl;
    return 0;
}
\end{consolethree}

\begin{ex}
Rewrite the following using enumerators:

\begin{consolethree}[escapeinside=||]
const int FEE = 1;
const int FIE = 0;
const int FOE = 2;
const int FUM = 3;
\end{consolethree}

Omit the enumeration name and integer values.
\end{ex}


\newpage\EMPHASIZE{Symbolic constants}

Sometimes we don't really care about the integer value
of enumerators. We really want the names of the enumerators to make it
easier to read our programs. Look at this example:

\begin{consolethree}[escapeinside=||]
enum GameState
{
    INIT, MENU, STARTING, RUN, RESTART, EXIT
};

GameState gameState = INIT;

int main()
{
    int energy = 1000;

    while (gameState != EXIT)
    {
        ...
        if (energy == 0)
            gameState = EXIT
        ...

        switch (gameState)
        {
            case INIT: ...
            case MENU: ...
            case STARTING: ...
            case RUN: ...
            case RESTART: ...
            case EXIT: ...
        }
    }

    return 0;
}
\end{consolethree}

Notice that it doesn't matter if the integer value of
\textit{INIT} is 0 or 1 or 2 or what-have-you.

Using this technique of programming where the code to execute is based
on the value of a variable, you create something called a finite state
machine.

\begin{ex}
Given

\begin{consolethree}[escapeinside=||]
enum EnumA {x, y, z};
enum EnumB {a, b, c};
\end{consolethree}

does this work?

\begin{consolethree}[escapeinside=||]
EnumA i = a;
\end{consolethree}
\end{ex}

\begin{ex}
Given

\begin{consolethree}[escapeinside=||]
enum EnumA {x, y, z};
enum EnumB (a, b, c};
\end{consolethree}

does this work?

\begin{consolethree}[escapeinside=||]
EnumA a0 = x;
EnumB a1 = a;
std::cout << (a0 == a1) << std::endl;
\end{consolethree}
\end{ex}


\newpage\EMPHASIZE{Multi-file Compilation}

Enumerations can be placed in a header file if other files (source or
header) use it.


\newpage\EMPHASIZE{Summary}

You can create symbolic constants using enum.

You can give the set of constants a name. This creates a type.

You can then create a variable of an enum type.

\begin{consolethree}[escapeinside=||]
enum PlaneState { STATIONARY, FLYING, MAINTENANCE };
PlaneState boeing747a = STATIONARY;
PlaneState boeing747b = FLYING;
\end{consolethree}

By default enumerators are given integer values: 0, 1, 2, ...

You can give the enumerators values:

\begin{consolethree}[escapeinside=||]
enum PlaneState
{
    STATIONARY=5, FLYING, MAINTENANCE=10
}; // int(FLYING) is 6
\end{consolethree}

The name of the enumeration is optional.

\end{document}